% !PRINTED 109
\weberpage{109}
The values of $\nu$ occurring in (20) can be arranged in pairs as follows:
\[
\nu,\ n-\nu, \qquad \nu = 1,\ 2,\ \cdots\ \frac{n-1}{2},
\]
and since
\[
\begin{aligned}
\vartheta_{00}\left(\frac{\nu}{n}\right) &= \vartheta_{00}\left(\frac{n-\nu}{n}\right), \quad \vartheta_{10}\left(\frac{\nu}{n}\right) = -\vartheta_{10}\left(\frac{n-\nu}{n}\right), \\
\vartheta_{01}\left(\frac{\nu}{n}\right) &= \vartheta_{01}\left(\frac{n-\nu}{n}\right),
\end{aligned}
\]
(20) agrees, apart from the sign, with the square of

\weq{21}{\dfrac{\displaystyle\prod_{1,\,\frac{n-1}{2}}^{\nu} \vartheta_{00}\left(\dfrac{\nu}{n}\right)\vartheta_{10}\left(\dfrac{\nu}{n}\right)\vartheta_{01}\left(\dfrac{\nu}{n}\right)}{\vartheta_{00}^{\frac{n-1}{2}}\vartheta_{10}^{\frac{n-1}{2}}\vartheta_{01}^{\frac{n-1}{2}}}.}

The last quotient, which is a continuous function of $\omega$ different from zero as long as the imaginary part of $\omega$ is positive, therefore likewise remains unchanged when $\omega$ is replaced by $2\omega$, and hence also by $4\omega$, $8\omega$, $\ldots$. The value of this expression can be determined by taking the imaginary part of $\omega$ to be infinite, and hence assuming $q=0$. But for $q=0$ we have (by § 25):
\[
\vartheta_{00}(v) = 1, \quad \vartheta_{01}(v) = 1, \quad \frac{\vartheta_{10}(v)}{\vartheta_{10}} = \cos\pi v,
\]
and hence the value of (21) is:

\weq{22}{\prod_{1,\,\frac{n-1}{2}}^{\nu} \cos\frac{\nu\pi}{n}.}

Since here $\nu\pi/n$ is less than $\frac{1}{2}\pi$, this product has a positive value. But we have
\[
\cos\frac{\nu\pi}{n} = -\cos\frac{(n-\nu)\pi}{n}
\]
and by Bd. I, § 144 we have
\[
2^{\frac{n-1}{2}} \prod_{1,\,\frac{n-1}{2}}^{\nu} \cos\frac{\nu\pi}{n} = 1.
\]
This proves the formula:

\weq{23}{2^{\frac{n-1}{2}} \prod_{1,\,\frac{n-1}{2}}^{\nu} \vartheta_{00}\left(\frac{\nu}{n}\right)\vartheta_{10}\left(\frac{\nu}{n}\right)\vartheta_{01}\left(\frac{\nu}{n}\right) = \vartheta_{00}^{\frac{n-1}{2}}\vartheta_{10}^{\frac{n-1}{2}}\vartheta_{01}^{\frac{n-1}{2}}.}
